\section{Evaluation Design}

\subsection{Exact release identity}

The evaluation is reconstructed from the successful v5.0.1 release-evidence workflow, run
\EvidenceRunRaw{} attempt \EvidenceRunAttemptRaw{}, whose head is the peeled tag commit
\IdentityText{\ProductCommitRaw}. The annotated tag object is
\IdentityText{\ProductTagObjectRaw}. The primary whitepaper artifact is
\texttt{wp-intake-bundle-v4-1}, artifact ID \EvidenceArtifactIdRaw{}. Its inner ZIP SHA-256 is
\IdentityText{\EvidenceBundleShaRaw}; its attested GitHub artifact digest and the contract digest
are retained in the companion identity file.

The main characterization run UUID is \IdentityText{\RunId}. The canonical generated-fixture hash
is \IdentityText{\DatasetHash}. The runtime image digest is
\IdentityText{\RuntimeImageDigestRaw}. These values are checked by build tests and by the companion
verifier rather than copied only into prose.

\subsection{Primary fixture}

The generated credit-decision fixture has eight fields: five numeric predictors, gender, race, and
the binary decision label. The historical slice has 3,000 rows. Amplification and intrinsic each
have 10,000 generated rows. The gender comparison is female versus male. Race contains Asian,
black, Hispanic, other, and white categories as declared by this fixture; it is not a general
ontology and does not cover all legally protected bases or intersectional combinations.

The evaluation questions are deliberately separated:

\begin{enumerate}
  \item Does the evidence path detect the fixture's known gender selection-rate disparity?
  \item Does amplification retain that signal at the configured output size?
  \item Does the intrinsic post-label control execute and remain separated from amplification?
  \item What distribution, correlation, demographic, exact-reproduction, and utility limitations
        accompany those fairness summaries?
  \item Do the configured Gold scenarios complete consistently across their fixed seed plan?
\end{enumerate}

\subsection{Evidence surfaces and priority}

When similarly named fields disagree, this paper uses an explicit source priority:

\begin{enumerate}
  \item \path{intake/metrics_uncertainty.json} and \path{intake/fairness_slices.json} for headline
        fairness values and methods;
  \item branch-specific quality certificates for branch quality and privacy-scope fields;
  \item the exact Gold index and the path-free merged-summary projection for multi-seed
    characterization;
  \item the paper-owned utility summary for the additional TSTR study; and
  \item the paper-owned regulatory mapping dated 5 August 2026 for current legal and governance
        context.
\end{enumerate}

The producer's original \path{intake/regulatory_matrix.csv} remains byte-preserved to maintain
provenance, but its narrative mapping is superseded for this publication by the dated paper-owned
overlay. Interpretive corrections do not alter any producer metric, interval, count, certificate,
or hash.

\subsection{Display and disclosure rules}

Gender is the headline protected-attribute result because the fixture and scenario are designed
around that positive control. Race remains machine-readable but is not plotted in the main result
figures because the Hispanic group has \RaceMinimumGroupN{} of 10,000 rows
(\RaceMinimumGroupPct{}\%). It \RaceCountFloorRelation{} the configured
\RaceDisplayMinGroupN{}-row count condition and \RaceShareFloorRelation{} the configured
\RaceDisplayMinGroupPct{}\% share condition. The paper still discloses the configured and effective
reference groups, selection policy, worst pair, minimum group count, and reasons for suppression.
Suppression is a presentation rule, not a claim that the race result is favorable, immaterial, or
absent.
